\DocumentMetadata{
  pdfversion=2.0,
  pdfstandard=ua-2,
  lang=en-US,
  tagging=on,
  tagging-setup={math/setup=mathml-SE,tabsorder=structure}
}
\documentclass[11pt,letterpaper]{article}
\usepackage[margin=0.68in,includefoot]{geometry}
\usepackage{newcomputermodern}
\usepackage{amsmath,amscd,mathtools}
\providecommand{\square}{\mdlgwhtsquare}
\usepackage[hidelinks]{hyperref}
\hypersetup{
  pdftitle={Logic PhD Exam, January 2023},
  pdfauthor={University of Florida Department of Mathematics},
  pdfsubject={Graduate mathematics examination},
  pdfdisplaydoctitle=true,
  bookmarksopen=true
}
\setcounter{secnumdepth}{0}
\setlength{\parindent}{0pt}
\setlength{\parskip}{0.28em}
\setlength{\emergencystretch}{3em}
\setlength{\leftmargini}{2.15em}
\setlength{\leftmarginii}{2.65em}
\setlength{\leftmarginiii}{3em}
\renewcommand{\labelenumi}{\textbf{\theenumi.}}
\pagestyle{empty}
\raggedbottom
\allowdisplaybreaks
\newenvironment{examproblems}[1]{%
  \begin{enumerate}%
  \setcounter{enumi}{#1}%
  \setlength{\topsep}{0.35em}%
  \setlength{\partopsep}{0pt}%
  \setlength{\itemsep}{0.48em}%
  \setlength{\parsep}{0.18em}%
}{\end{enumerate}}
\newenvironment{examsubparts}{%
  \begin{enumerate}%
  \setlength{\topsep}{0.18em}%
  \setlength{\partopsep}{0pt}%
  \setlength{\itemsep}{0.16em}%
  \setlength{\parsep}{0.1em}%
}{\end{enumerate}}
\newenvironment{examinstructions}{%
  \par\begingroup\itshape
}{%
  \par\endgroup\vspace{0.18em}
}
\makeatletter
\renewcommand\section{\@startsection{section}{1}{0pt}%
  {-1.25ex plus -.25ex minus -.1ex}%
  {0.55ex plus .1ex}%
  {\normalfont\large\bfseries}}
\renewcommand{\@maketitle}{%
  \newpage\null\vskip -1.2em%
  {\footnotesize\bfseries
    \href{https://www.ufl.edu/}{University of Florida}, \href{https://math.ufl.edu/}{Department of Mathematics}\par}%
  \vskip 0.18em%
  \tagstructbegin{tag=Title}%
    {\LARGE\bfseries\href{https://gma.math.ufl.edu/exams/logic-phd/}{Logic PhD Exam}, January 2023\par}%
  \tagstructend%
  \vskip 0.35em%
  \tagmcbegin{artifact}%
    \hrule height 1pt%
  \tagmcend%
  \vskip 0.55em%
}
\makeatother
\title{Logic PhD Exam, January 2023}
\author{}
\date{}
\begin{document}
\maketitle
\thispagestyle{empty}
\begin{examinstructions}
Solve 5 problems of the following; at least one from each section.
\end{examinstructions}
\section{A. Set Theory.}
\begin{examproblems}{0}
\item
Let~\(X\) be an uncountable set. Let \(f:X\to[X]^{<\aleph_0}\) be a function. Show that there are distinct points \(x_0,x_1\in X\) such that \(x_0\notin f(x_1)\) and \(x_1\notin f(x_0)\).
\item
Let~\(P\) be a forcing and~\(\kappa\) be a regular cardinal larger than \(|P|\). Show that \(P\Vdash\check{\kappa}\) is a cardinal.
\item
What is the perfect set theorem in descriptive set theory? State it and prove it.
\end{examproblems}
\section{B. Computability.}
\begin{examproblems}{0}
\item
Choose one of the standard definitions of a recursive function, and show that the increasing enumeration of the Fibonacci numbers (defined by \(f(0)=0\), \(f(1)=1\) and \(f(n+2)=f(n)+f(n+1)\)) is a recursive function.
\item
Provide an example of sets \(A,B\subseteq\mathbb{N}\) such that~\(A\) is Turing reducible to~\(B\) but not many-one reducible to the complement of~\(B\).
\item
Prove that a simple c.e. set exists, where a set \(A\subseteq\mathbb{N}\) is simple if the complement of~\(A\) has no infinite c.e. subset.
\end{examproblems}
\section{C. Model theory.}
\begin{examproblems}{0}
\item
Let~\(\mathcal{L}\) be a language containing two binary relational symbols \(E_0,E_1\). Let~\(T\) be the theory stating that \(E_0,E_1\) are equivalence relations and \(E_0\subseteq E_1\). Let~\(\mathcal{F}\) be the class of finite models of~\(T\). Show that~\(\mathcal{F}\) is a Fraïssé class and provide a simple description of its limit.
\item
State the Łoś theorem and prove it.
\item
Find a complete theory with more than one countably infinite model up to isomorphism. Provide two nonisomorphic countable models of the theory.
\end{examproblems}
\end{document}
